\scanpage{56}

Algèbres étales et la différente (2)\hfill Tribu n° 65

\uline{n° 2.}\quad \uline{Algèbres étales et algèbres non ramifiées.}

Soit $B$ une $A$-algèbre. On considère les propriétés suivantes :

a)\quad $B$ est de p. f.

b)\quad $B$ est plate

c) Pour tout $q\in\operatorname{Spec}B$ et pour $p=q\cap A$, $B_q/pB_q$ est une extension séparable finie de $A_p/pA_p$.

En présence de a) \deleted{et b)}, c) veut dire que $\Omega_{B/A}$ est nul en $q$.

On dit que $B$ est étale si on a a), b), c).

On dit que $B$ est étale en $q\in\operatorname{Spec}(B)$ s'il existe $f\notin q$ tel que on ait a), b), c) pour $B_f$. (i. e. tel que $B_f$ soit étale).

On dit que $B$ est non ramifiée si elle est de p. f. et si $\Omega_{B/A}=0$.

Ainsi : étale = non ramifiée + plate.

Si $A$ est réduite, $B$ est réduite.

Transitivité ($A\longrightarrow B\longrightarrow C$) de l'étalité et de la non-ramification.

Effet d'un changement d'anneau sur ces notions : la montée marche dans le cas général, la descente marche dans le cas fidèlement plat.

\uline{n° 3.}\quad \uline{Algèbres finies étales et non ramifiées.}

Les équivalences de la prop. 1, p. 1, du texte (améliorées comme dans la Tribu n° 64 ; on supposera seulement $B$ de présentation finie sur $A$).

Soient $B$ de p. f. sur $A$ noethérien, $q\in\operatorname{Spec}B$ et $p=q\cap A$ ; pour que $B$ soit étale en $q$, il faut et il suffit que $\hat{B}_q$ soit étale finie sur $\hat{A}_p$ ; analogue non ramifié ; la démonstration utilise la platitude supérieure. Corollaire pour le cas où les corps résiduels sont égaux (cf. S. G. A.).

Pour que $B$ soit étale finie sur $A$ connexe, il faut et il suffit qu'il existe $A'$ fidèlement plate sur $A$ telle que $A'\otimes B$ soit un produit de copies de $A'$ (si $A$ n'est pas connexe, ça ne marche que par morceaux).

La prop. 5 du texte (arrangée comme dans la Tribu n° 64). Le calcul de traces de la prop. 3, a) du § 2 (p. 16).

\uline{n° 4.}\quad \uline{Structure locale des algèbres non ramifiées et étales.}

Soit $B$ une $A$-algèbre finie et soit $q\in\operatorname{Spec}(B)$. Pour que $B$ soit non ramifiée en $q$ il faut et il suffit qu'il existe $x\in\unsure{B'|B}$ de polynôme minimal $F$

\mspage{55}
